\chapter{Timing Analysis and Parasitics}
\label{ch:timing}

\glance{OpenROAD's timing engine is OpenSTA, adapted to OpenDB by
\pkg{src/dbSta}. Before routing, wire RC is \emph{estimated} from Steiner trees
(\pkg{src/stt} + \pkg{src/est}); after detailed routing it is \emph{extracted}
by OpenRCX (\pkg{src/rcx}) and back-annotated into STA. Every optimizer ---
especially the resizer --- asks STA for slack and uses these parasitics.}

\section{OpenSTA and the database bridge}

OpenSTA (the Parallax/UCSD static timing analyzer) is vendored under
\pkg{src/sta}. OpenROAD does not reimplement STA; it adapts OpenSTA to OpenDB
through \pkg{src/dbSta}:

\begin{itemize}
  \item \id{sta::dbNetwork} maps OpenDB objects to STA's abstract netlist
        (\id{dbInst}$\leftrightarrow$\id{Instance},
        \id{dbNet}$\leftrightarrow$\id{Net},
        \id{dbITerm}/\id{dbBTerm}$\leftrightarrow$\id{Pin}).
  \item \id{sta::dbSta} is the OpenSTA engine wired to that network; it is
        constructed once in \id{OpenRoad::init()} and shared by every module.
\end{itemize}

OpenSTA owns the timing graph, delay calculation, path search, and SDC
handling. Its user commands are the familiar ones:

\begin{center}
\small
\begin{tabularx}{\linewidth}{@{}lX@{}}
\toprule
\textbf{Command} & \textbf{Purpose} \\
\midrule
\cmd{read\_liberty} & load a Liberty (.lib) timing library \\
\cmd{read\_sdc} / \cmd{create\_clock} / \cmd{set\_input\_delay} & timing constraints \\
\cmd{report\_checks} & path reports (setup/hold) \\
\cmd{report\_wns} / \cmd{report\_tns} & worst / total negative slack \\
\cmd{report\_power} & static/dynamic power estimate \\
\bottomrule
\end{tabularx}
\end{center}

Engine internals are out of scope here --- see the standalone OpenSTA docs.

\section{Parasitic estimation (pre-route): \pkg{src/est} + \pkg{src/stt}}

Before wires exist, timing still needs RC. OpenROAD builds an estimated
topology and applies per-unit R/C:

\begin{enumerate}
  \item \pkg{src/stt} (\id{stt::SteinerTreeBuilder}) builds a
        \textbf{rectilinear Steiner tree} for each net from pin locations
        (FLUTE / Prim--Dijkstra). \cmd{set\_routing\_alpha} trades wirelength
        ($\alpha{=}0$) against path depth ($\alpha{=}1$); typical values
        $0.4$--$0.8$.
  \item \pkg{src/est} (\id{est::EstimateParasitics}) annotates that topology
        with R/C from \cmd{set\_wire\_rc} (or from a LEF layer via
        \cmd{-layer}).
\end{enumerate}

\begin{lstlisting}[style=ortcl]
set_wire_rc -signal -layer metal3
set_wire_rc -clock  -layer metal5
estimate_parasitics -placement          ;# after placement
# later, after global_route:
estimate_parasitics -global_routing     ;# use GR topology/layers
\end{lstlisting}

\cmd{set\_layer\_rc} can override LEF R/C per layer or via;
\cmd{report\_layer\_rc} dumps the values in use.

\section{Parasitic extraction (post-route): \pkg{src/rcx}}

After detailed routing, OpenRCX (\id{rcx::Ext}) extracts real parasitics from
the geometry: coupling capacitance as a function of neighbour distance and
over/under context, plus wire and via resistance. A once-per-process
\emph{extraction rules} file (RC technology) is loaded with
\cmd{set\_extraction\_rules\_file} / \cmd{define\_process\_corner}; then:

\begin{lstlisting}[style=ortcl]
define_process_corner -ext_model_index {0 typical} rcx_patterns.rules
extract_parasitics
write_spef design.spef          ;# optional SPEF dump
\end{lstlisting}

Extracted R/C is stored on OpenDB objects with pointers back to the wires and
vias, and is visible to STA for signoff-quality timing.

\section{Summary}

\begin{table}[h]
\centering
\small
\begin{tabularx}{\linewidth}{@{}lllX@{}}
\toprule
\textbf{Task} & \textbf{Module} & \textbf{Key command} & \textbf{Class} \\
\midrule
STA engine         & \pkg{src/sta}+\pkg{dbSta} & \cmd{report\_checks} & \id{sta::dbSta} \\
Steiner topology   & \pkg{src/stt} & \cmd{set\_routing\_alpha} & \id{stt::SteinerTreeBuilder} \\
RC estimation      & \pkg{src/est} & \cmd{estimate\_parasitics} & \id{est::EstimateParasitics} \\
RC extraction      & \pkg{src/rcx} & \cmd{extract\_parasitics} & \id{rcx::Ext} \\
\bottomrule
\end{tabularx}
\end{table}

\glance{Rule of thumb: \emph{before} route use \cmd{estimate\_parasitics};
\emph{after} detailed route use \cmd{extract\_parasitics}. The resizer
(Ch.~\ref{ch:rsz}) consumes whichever is currently annotated.}
